\section*{Ethics Statement}
\label{sec:ethics}

\paragraph{Data and licences.}
We use public data under their licences: Natural Questions and NQ-open (CC BY-SA 3.0) through the BEIR distribution \citep{thakur2021beir}.
The questions are anonymised search queries over Wikipedia; we add no personal data and collect none.
All models are openly licensed: all-MiniLM-L6-v2, Qwen3-Embedding-0.6B, Qwen2.5-7B-Instruct-AWQ and Mistral-7B-Instruct-v0.3 under Apache-2.0, and the bge and e5 encoders under MIT.
We use each dataset and model for research evaluation, which is consistent with its intended use.
We release our derived data under CC BY-SA and our code under an open licence.

\paragraph{Risks.}
The study evaluates existing retrieval methods and creates no new model or capability.
A reader could over-generalise our comparison of models to other readers, corpora or languages; we scope every claim to our setup and list its limits.
Faithfulness scores come from an unvalidated automatic judge and should not be used as evidence that a system is safe to deploy.

\paragraph{Human annotation.}
One annotator from the author team labels the \NUM{n.human}-item validation sample with the same instructions the judge receives (Appendix~\ref{app:judge}).
No external annotators or participants are involved, so recruitment, payment and ethics review do not apply.

\paragraph{Compute and use of AI assistants.}
All experiments ran on free Kaggle sessions (Tesla T4); the study had no monetary cost.
We used an AI assistant (Claude) to help write analysis and plotting code and to draft and edit text.
We designed the study, ran all experiments, checked every reported number against the released result files and take full responsibility for the content.
